\begin{textsample}{NEG Dim 5 – pb – Score -1.00 – pb/pb\_ef\_53.txt}  \label{ex:f5_neg_193}
9.1 Ensino Religioso 9.1.1 Introdução O Ensino Religioso é um componente curricular presente na BNCC essencial e determinante na formação do ser humano.
Todos os componentes curriculares são importantes, mas o Ensino Religioso contém o elemento religioso que alicerça a vida humana.
Segundo Coll ( 2002 ), a sociedade é resultado, não só do logos ( racionalidade ), mas, sobretudo, das dimensões mítico-simbólica e do mistério.
A primeira corresponde a um nível de realidade da razão lógica, reflexiva e conceitual.
Já o mistério se relaciona ao que não pode ser pensado e nem dito, mas que igualmente real, tal qual percebemos pela razão.
Assim, toda cultura veicula uma concepção de humano, do divino e do cósmico. ( POZZER et al., 2015 ).
Ao longo da nossa história, inúmeras maneiras foram construídas para perpetuar o analfabetismo religioso, como objetivo de desvalorizar e negar a diversidade de crenças e filosofia de vidas, sobretudo das comunidades e povos subjugados, marginalizados e inferiorizados.
De acordo com a BNCC, > O Ensino Religioso assumiu diferentes perspectivas teórico-metodológicas, geralmente de viés confessional ou interconfessional.
A partir da década de 1980, as transformações socioculturais que provocaram mudanças paradigmáticas no campo educacional também impactaram o Ensino Religioso.
Em função dos promulgados ideais de democracia, inclusão social e educação integral, vários setores da sociedade civil passaram a reivindicar a abordagem do conhecimento religioso e o reconhecimento da diversidade religiosa no âmbito dos currículos escolares.
A temática do Ensino Religioso Não Confessional nas redes públicas de ensino como direito do educando e como obrigatoriedade de oferta por parte do Estado está intrinsecamente ligada aos princípios democráticos e da paz, aos direitos civis e políticos de cada cidadão, bem como dos Direitos Humanos.
O Brasil confirma no Art. 5º de sua Carta Magna que “ Todos são iguais perante a lei, sem distinção de qualquer natureza, garantindo -se aos brasileiros e aos estrangeiros residentes no País a inviolabilidade do direito à vida, à liberdade, à igualdade, à segurança e à propriedade. ” Nas duas últimas décadas, várias instituições assumiram a vanguarda em prol das lutas e conquistas do Ensino Religioso Não Confessional, com base na Constituição Federal de 1988 ( artigo 210 ) e na LDB nº 9.394/1996 ( artigo 33, alterado pela Lei nº 9.475/1997 ).
A legislação estabeleceu os princípios e os fundamentos que, conforme normatiza a BNCC, devem alicerçar epistemologias e pedagogias do Ensino Religioso, cuja função educacional, enquanto parte integrante da formação básica do cidadão, é assegurar o respeito à diversidade cultural religiosa, sem proselitismos.
Mais tarde, a Resolução CNE/CEB nº 04/2010 e a Resolução CNE/CEB nº 07/2010 reconheceram o Ensino Religioso como uma das cinco áreas de conhecimento do Ensino Fundamental de 09 ( nove ) anos.[1 ] Os Parâmetros Curriculares Nacionais do Ensino Religioso (PCNER)[2 ] ressaltam a construção de culturas de paz, o respeito ao outro e às relações socioambientais, o direito à vida do/no Planeta, a eliminação das diversas formas de preconceitos e \textbf{intolerâncias}, a educação intercultural, os direitos humanos e a aprendizagem, a laicidade do Estado, a diversidade e liberdade religiosa, pessoas sem religião e outras temáticas que integram o currículo da escola na atualidade.
Uma vez que o Ensino Religioso foi estabelecido como componente curricular de oferta obrigatória nas escolas públicas de Ensino Fundamental, com matrícula facultativa para o aluno, não apenas na Paraíba, mas também por todo o país, teve início o processo de elaboração de materiais didático-pedagógicos, propostas curriculares, além de cursos de formação inicial e continuada “ que contribuíram para a construção da área do Ensino Religioso, cujas natureza e finalidades pedagógicas são distintas da confessionalidade ” ( BNCC ).
Considerando os marcos normativos e, em conformidade com as competências gerais estabelecidas no âmbito da BNCC, o Ensino Religioso deve atender aos seguintes objetivos : a ) Proporcionar a aprendizagem dos conhecimentos religiosos, culturais e estéticos, a partir das manifestações religiosas percebidas na realidade dos educandos ; b ) Propiciar conhecimentos sobre o direito à liberdade de consciência e de crença, no constante propósito de promoção dos direitos humanos ; c ) Desenvolver competências e habilidades que contribuam para o diálogo entre perspectivas religiosas e seculares de vida, exercitando o respeito à liberdade de concepções e o pluralismo de ideias, de acordo com a Constituição Federal ; d ) Contribuir para que os educandos construam seus sentidos pessoais de vida a partir de valores, princípios éticos e da cidadania.
Esses objetivos tecem uma teia de saberes, num investimento para as gerações presentes e futuras, reconhecendo a existência de relacionamentos humanos essenciais aos valores da vida, num processo de humanização da educação, através do cuidar e do educar.
A respeito do conhecimento religioso, objeto da área de Ensino Religioso, a BNCC assegura que > é produzido no âmbito das diferentes áreas do conhecimento científico das Ciências Humanas e Sociais, notadamente da(s) Ciência(s) da(s) Religião(ões ).
Essas Ciências investigam a manifestação dos fenômenos religiosos em diferentes culturas e sociedades, enquanto um dos bens simbólicos resultantes da busca humana por respostas aos enigmas do mundo, da vida e da morte.
De modo singular, complexo e diverso, esses fenômenos alicerçaram distintos sentidos e significados de vida e diversas ideias de divindade(s), em torno dos quais se organizaram cosmovisões, linguagens, saberes, crenças, mitologias, narrativas, textos, símbolos, ritos, doutrinas, tradições, movimentos, práticas e princípios éticos e morais.
Os fenômenos religiosos em suas múltiplas manifestações são parte integrante do substrato cultural da humanidade.
É sempre importante ressaltar que o Ensino Religioso deve tratar os conhecimentos religiosos a partir das diversas culturas e tradições religiosas, baseando -se em pressupostos éticos e científicos, sem privilégio de nenhuma crença ou convicção.
Por esse motivo, a BNCC reafirma a interculturalidade e a ética da alteridade como “ fundamentos teóricos e pedagógicos do Ensino Religioso, porque favorecem o reconhecimento e respeito às histórias, memórias, crenças, convicções e valores de diferentes culturas, tradições religiosas e filosofias seculares de vida. ” Com base nisso, é possível afirmar que é o objetivo principal de Ensino Religioso é, em conformidade com a BNCC, “ construir, por meio do estudo dos conhecimentos religiosos e das filosofias de vida, atitudes de reconhecimento e respeito às alteridades ”.
Assim, conclui -se que tal componente curricular se constitui em “ um espaço de aprendizagens, experiências pedagógicas, intercâmbios e diálogos permanentes, que visam o acolhimento das identidades culturais, religiosas ou não, na perspectiva da interculturalidade, direitos humanos e cultura da paz ” ( BNCC ). [ 1 ] : BRASIL.
Constituição da República Federativa do Brasil ( 1988 ).
Brasília, DF : Senado Federal, 1988.
Disponível em :.
BRASIL.
Lei nº 9.394, de 20 de dezembro de 1996.
Estabelece as diretrizes e bases da educação nacional.
Diário Oficial da União, Brasília, 23 de dezembro de 1996.
Disponível em :.
BRASIL.
Conselho Nacional de Educação / Câmara de Educação Básica.
Resolução nº 4, de 13 de julho de 2010.
Define Diretrizes Curriculares Nacionais Gerais para a Educação Básica.
Diário Oficial da União, Brasília, 14 de julho de 2010, Seção 1, p. 824.
Disponível em :.
BRASIL.
Ministério da Educação.
Conselho Nacional de Educação ; Câmara de Educação Básica.
Resolução nº 7, de 14 de dezembro de 2010.
Fixa Diretrizes Curriculares Nacionais para o Ensino Fundamental de 9 ( nove ) anos.
Diário Oficial da União, Brasília, 15 de dezembro de 2010, Seção 1, p. 34.
Disponível em :.
Acesso em : 7 nov. 2017. [ 2 ] : Parâmetros Curriculares Nacionais do Ensino Religioso apresentam os 5 eixos temáticos do ER, organizam os conteúdos respeitando à diversidade cultural religiosa existente no contexto escolar. ( FONAPER, 2009 ).

% matched lemmas: intolerância
\end{textsample}
